\documentclass{article}
\usepackage{graphicx}
\usepackage{epsfig}
\usepackage{amsmath}
\usepackage{amsfonts}
\usepackage{amssymb}
\usepackage{latexsym}
\usepackage{mathrsfs}
\usepackage{algpseudocode}

\oddsidemargin -0.0 in \topmargin -25 mm \textheight 10in
\textwidth 6.5 in

%\pagenumbering{empty}
\begin{document}
\begin{centering}

\textbf{\large Assignment 5 CSE 16-02: Applied Discrete Mathematics}\\


\textbf{\large Winter 2026}\\
Instructor: Cedric Westphal  \\
Due February 12th, 2026 11:59 PM\\

\end{centering}

\vspace{12pt}
\noindent Complete the questions carefully and show your work. Please remember to:
\begin{itemize}
	\item Scan the pages and combine them into a single PDF file. Make sure your work is readable and that pages are oriented correctly.
	\item Upload your PDF file to Canvas/Gradescope under Assignment 5.
	\item In Gradescope select the page(s) for each problem.
\end{itemize}

\noindent Notation: $\emptyset$ is the empty set, $\mathbb{N}$ is the set of natural numbers, $\mathbb{Z}$ is the set of integers, $\mathbb{Q}$ is the set of rational numbers, and $\mathbb{R}$ is the set of real numbers. $\mathscr{P}(A)$ is the power set of $A$.\\
\textbf{[46 points total in assignment.]}

\begin{enumerate}
\setcounter{enumi}{-1} 
\item Each answer is clearly written and each answer is correctly mapped to its rubric on Gradescope. [4 points]

\item  For $a,b \in \mathbb{N}$, prove the following statements below. Hint: to prove that $g$ is the gcd($a,b$) you need to prove that $g | a$ and $g|b$ and that $g$ is the greatest such common divisor.\textbf{[10 points]}
\begin{enumerate}
    \item Let $a = 5x$ and $b = 5y+3$ for some $x$ and $y$ in $\mathbb{N}$. Prove that gcd($a,b$) = gcd($x,b$). Write a proof that shows: (i) that gcd($a,b)| x$, then (ii) show that it implies that gcd($a,b$) equals gcd($x,b$). \textbf{[2 points]}
    \item If both $a$ and $b$ are divisible by 3, then gcd($a,b$) = 3 gcd($a/3,b/3$). \textbf{[2 points]}
    \item If $m$ is a common divisor of $a$ and $b$, show the general case of the previous question, namely that gcd($a,b$) = $m$ gcd($a/m,b/m$). \textbf{[2 points]}
    \item Prove that gcd($a,b$) = gcd($a+b,b$). \textbf{[2 points]}
    \item Is it true or false that, for $a>b$, gcd($a,b$) = gcd($a+b,a-b$)? Give a clear explanation of your reasoning. \textbf{[2 points]}
\end{enumerate}

\item We will prove that if $n$ is an odd integer $\in \mathbb{N}$, then 8 divides $n^6 - 1$. We're going to do this in two ways. \textbf{[7 points]} 
\begin{enumerate}
    \item First, prove that if $n$ is an odd integer $\in \mathbb{N}$, then 8 divides $n^2 - 1$.
    \item Prove that if $n$ is an odd integer, then so is $n^3$.
    \item Prove that 8 divides $n^6-1$ for $n$ odd, using the previous questions.
    \item For the second way to prove this, express $n^6-1$ for $n$ odd using the binomial theorem for $n = 2k+1$ as a sum of terms of the form $a_i k^i$? Write this sum.
    \item Consider the last two terms for $a_2 k^2$ and $a_1 k$. Can you rewrite $a_2 k^2 + a_1 k$ as $b k^2 + c (k^2 + k)$ where $b$ is divisible by 8?
    \item Prove that 8 divides $c (k^2 +k)$
    \item Prove the result stated at the beginning of this Problem, using the previous three steps. 
\end{enumerate}

 \item Prove or disprove:  $\forall x \in \mathbb{R}, \sqrt{x} \notin \mathbb{Q} \rightarrow x \notin \mathbb{Q} $.   \textbf{[5 points]}

\item Prove or disprove: $ \sqrt{3} \in \mathbb{Q}$.  \textbf{[5 points]} % $\forall x \in \mathbb{R}, x \notin \mathbb{Q} \rightarrow \sqrt[3]{x} \notin \mathbb{Q} $.  \textbf{[5 points]}

\item Prove:  If $n, m \in \mathbb{Z}$ and $7 n^2 + 11 m^3 = 325,014$ then $n$ and $m$ have the same parity.  \textbf{[5 points]}

\item Prove or disprove: If $x \notin \mathbb{Q}$, $y \in \mathbb{Q}$, $x \neq 0$, and $y \neq 0$, then for all $n, m \in \mathbb{N}$,  $x^ny^m
%{\frac{1}{m}}  
\notin \mathbb{N}$.  \textbf{[5 points]}

\item Consider the sets 
\[ S_1 = \{k^2+3: k \mbox{ is an even natural number} \},   S_2 = \{4k+3: k \in \mathbb{N} \}.    \]
\begin{enumerate}
	\item Prove $S_1 \subseteq S_2$. \textbf{[3 points]}
	\item Prove $\neg (S_2 \subseteq S_1)$. \textbf{[2 points]}
\end{enumerate}
\end{enumerate}

\end{document} 
